\chapter{Kéo Cả Lớp Đứng Dậy}
\chapterintro{Đổi tư thế để đổi nhịp}{Có những lớp chỉ cần rời ghế một phút là tiết học sống lại rõ rệt.}

\section{Đứng lên đổi chỗ mười giây}
\begin{khoidongbox}{Hoạt động khởi động}
Cho lớp đứng lên, bước đổi chỗ thật nhỏ với bạn bên cạnh rồi ngồi lại ngay.
\end{khoidongbox}
\begin{visaohieuqua}
Sự thay đổi vị trí ngắn giúp cơ thể tỉnh lên mà không làm lớp loạn. Nó rất hợp với đầu giờ trì trệ.
\end{visaohieuqua}
\begin{cachlam5phut}
Quy định rõ: đổi chỗ mười giây, ngồi xuống, rồi nói một câu với bạn mới ngồi cạnh.
\end{cachlam5phut}
\ngancach

\section{Đứng theo phe đồng ý hay phân vân}
\begin{khoidongbox}{Hoạt động khởi động}
Đưa một nhận định nhẹ và mời lớp đứng sang hai phía: \textit{đồng ý} hoặc \textit{còn phân vân}.
\end{khoidongbox}
\begin{visaohieuqua}
Đứng theo phe tạo không khí tham gia rất nhanh. Học sinh thấy mình được chọn vị trí trước khi phải nói.
\end{visaohieuqua}
\begin{cachlam5phut}
Chỉ dùng với câu nhận định không quá nhạy cảm. Sau đó mời mỗi phía đúng một tiếng nói đại diện.
\end{cachlam5phut}
\ngancach

\section{Bàn nào đứng chốt trước}
\begin{khoidongbox}{Hoạt động khởi động}
Chọn một bàn đứng lên trước và nói đúng một câu chốt về bài cũ hoặc chủ đề sắp học.
\end{khoidongbox}
\begin{visaohieuqua}
Khởi động theo bàn nhẹ áp lực hơn gọi cá nhân. Nó cũng làm lớp thấy hoạt động có tính đồng đội.
\end{visaohieuqua}
\begin{cachlam5phut}
Chỉ cần 2 bàn tham gia là đủ. Khen một điểm cụ thể rồi quay về nhịp chung của lớp.
\end{cachlam5phut}
\ngancach

\section{Đổi người nói khi đang đứng}
\begin{khoidongbox}{Hoạt động khởi động}
Khi một học sinh đang nói, cho em chọn ngay người tiếp theo để nối ý trước khi cả lớp ngồi xuống.
\end{khoidongbox}
\begin{visaohieuqua}
Cách này giữ nhịp lời nói liền mạch và khiến cả lớp chú ý hơn vì ai cũng có thể được trao lượt.
\end{visaohieuqua}
\begin{cachlam5phut}
Giới hạn đúng 3 lượt nối, mỗi lượt một câu. Không kéo dài thành hoạt động chính của tiết.
\end{cachlam5phut}
\ngancach

\section{Một vòng đứng cảm ơn ngắn}
\begin{khoidongbox}{Hoạt động khởi động}
Cho lớp đứng lên và mỗi dãy nói đúng một điều mình muốn cảm ơn ở buổi học trước hoặc trong lớp.
\end{khoidongbox}
\begin{visaohieuqua}
Đây là kiểu mở tiết vừa làm lớp ấm lại vừa không quá ồn. Nó đặc biệt hợp đầu tuần hoặc sau một buổi lớp căng.
\end{visaohieuqua}
\begin{cachlam5phut}
Mỗi dãy chỉ một câu. Sau đó mời lớp ngồi xuống và nói luôn mục tiêu hôm nay.
\end{cachlam5phut}